% !TEX root = ../main.tex

\chapter{Conclusions}

\section{Answers to the Research Questions}

\subsection{RQ1: Revision-Correct Representation}

The design answers RQ1 by separating logical paper identity from revision identity and carrying revision, content hash, transformation configuration, and lineage into derived artifacts.
This answer establishes a testable invariant.
The inspected source revision supports an implementation inventory, but the invariant still requires E1 execution logs before it can be established empirically.

\subsection{RQ2: Evidence-Bounded Assistance}

The design answers RQ2 by decomposing answers into claims and distinguishing identifier validity, evidence availability, consistency, and coverage.
It defines scoped retrieval, resolvable source spans, and abstention.
The literature supports the need for these distinctions~\cite{gao2023alce,es2024ragas,huang2024grounded,niu2024ragtruth}.
Whether the complete pipeline improves retrieval or citation quality remains \pending{} until E2 is executed.

\subsection{RQ3: Adaptive Interest Memory}

The design answers RQ3 with explicit interests, bounded implicit-event semantics, one bounded time-decayed normalized behavior-v1 vector, inspectable recommendation reasons, and user correction.
Separable short-term and long-term state remains planned behavior-v2 rather than a current implementation claim.
The behavior-v1 model's benefit for ranking, diversity, adaptation, or user control remains \pending{} until E3 provides temporal and correction evidence.

\subsection{RQ4: Bounded Mobile Integration}

The prototype demonstrates that the intended workflow can be represented as mobile interactions.
The supplied presentation reports that participants completed the digest, evidence, Q\&A, and preference tasks more often than the graph task.
Because raw records and a production build are absent, this evidence is \observed{}/\limited{}.
Mobile latency, reliability, privacy behavior, and redesign effectiveness remain unestablished.

\section{Contribution and Limitations}

The current contribution is an evidence-first system and evaluation framework.
It connects four mechanisms through explicit traceability and makes its own evidence state visible.
This framing avoids two common errors: describing a technology stack as novelty and treating fluent, citation-bearing output as verified understanding.

The same evidence-first standard limits the conclusion.
The inspected source snapshots support only the bounded implementation claims documented in this thesis; they do not justify claims of production readiness, improved RAG quality, successful behavioral adaptation, or general usability.
Those limitations do not invalidate the design contribution, but they prevent the thesis from claiming empirical completion.

\section{Prioritized Future Work}

The first priority is not an additional product feature.
It is an execution freeze at an inspected repository revision that permits a verified implementation inventory and reproducible evaluation runs.
The second priority is completion of E1--E3 with frozen data, baselines, configurations, outputs, and analyses.
The third priority is an E4 package containing system logs and participant-level usability records, followed by a focused retest of graph discoverability.

Only after the core evidence chain is complete should the project prioritize stretch features such as autonomous multi-hop research, voice input, camera OCR, collaboration, or on-device language models.
These features broaden capability but do not repair missing evidence for the central contribution.

\section{Final Statement}

MNEME is most defensible when understood not as an autonomous researcher, but as an accountable mobile loop for continuing literature work.
The loop is useful only if paper revisions remain identifiable, generated claims remain inspectable, behavioral adaptation remains correctable, and mobile integration preserves those properties.
This draft defines how each promise can be tested and marks where evidence is still required.

\clearpage
